\documentclass[10pt]{article}
% \usepackage[utf8]{inputenc}

% \usepackage[a4paper, total={6in, 8in}, margin = 3cm]{geometry}
% \usepackage[total={6in,8in}, margin = 3cm]{geometry}
\usepackage[a4paper, margin = 3cm]{geometry}

\usepackage[english]{babel}
\usepackage{authblk}
\usepackage{enumerate, enumitem}
\usepackage{setspace}
\usepackage{natbib}
% \usepackage{parskip}

%%% packages and libraries %%%
\usepackage{amsmath, amsthm, amssymb, amsfonts}
\usepackage{mathtools}
\usepackage{enumerate, enumitem}
\usepackage{algorithm, algpseudocode}
\usepackage{xparse}

%%% custom stats commands %%%
\newcommand{\simiid}{\overset{\scriptscriptstyle{\textrm{i.i.d.}}}{\sim}}
\newcommand{\indep}{\mathrel{\perp\!\!\!\perp}}
\newcommand{\estm}[2][]{\widehat{#2}_{#1}} % estimator
\newcommand{\iter}[2]{{#1}^{(#2)}} % iteration of algorithm

%%% custom math commands %%%
% \newcommand{\tpose}[1]{#1^{\scriptstyle{\mathsf{T}}}} % transpose
% \newcommand{\ctpose}[1]{#1^{\scriptstyle{\mathsf{H}}}} % conjugate transpose
\newcommand{\floor}[1]{\lfloor{#1}\rfloor}
\newcommand{\ceil}[1]{\lceil{#1}\rceil}

\DeclareMathOperator*{\argmax}{arg\,max}
\DeclareMathOperator*{\argmin}{arg\,min}
\DeclareMathOperator{\sign}{sign}
\DeclareMathOperator{\sgn}{sgn}
\DeclareMathOperator{\rank}{rank}
\DeclareMathOperator{\spn}{span}
\DeclareMathOperator{\tr}{tr}
\DeclareMathOperator{\supp}{supp}

% probability of event (optional subscript, blank by default)
\NewDocumentCommand{\prob}{ s O{} m }{
    \IfBooleanTF{#1}
        {\bbP_{#2}(#3)}              % \prob*{} (no resizing)
        {\bbP_{#2}\!\left(#3\right)} % \prob (auto-sizing)
}

% expectation (optional subscript, blank by default)
\NewDocumentCommand{\expec}{ s O{} m }{
    \IfBooleanTF{#1}
        {\bbE_{#2}[#3]}              % \expec*{} (no resizing)
        {\bbE_{#2}\!\left[#3\right]} % \expec (auto-sizing)
}

% variance (optional subscript, blank by default)
\NewDocumentCommand{\var}{ s O{} m }{
    \IfBooleanTF{#1}
        {\operatorname{Var}_{#2}(#3)}              % \var*{} (no resizing)
        {\operatorname{Var}_{#2}\!\left(#3\right)} % \var (auto-sizing)
}

% covariance (optional subscript, blank by default)
\NewDocumentCommand{\covar}{ s O{} m m }{
    \IfBooleanTF{#1}
        {\operatorname{Cov}_{#2}(#3, #4)}              % \covar*{} (no resizing)
        {\operatorname{Cov}_{#2}\!\left(#3, #4\right)} % \covar (auto-sizing)
}

% correlation (optional subscript, blank by default)
\NewDocumentCommand{\corr}{ s O{} m m }{
    \IfBooleanTF{#1}
        {\operatorname{corr}_{#2}(#3, #4)}              % \corr*{} (no resizing)
        {\operatorname{corr}_{#2}\!\left(#3, #4\right)} % \corr (auto-sizing)
}

% bias (optional subscript, blank by default)
\NewDocumentCommand{\bias}{ s O{} m }{
    \IfBooleanTF{#1}
        {\operatorname{bias}_{#2}(#3)}              % \bias*{} (no resizing)
        {\operatorname{bias}_{#2}\!\left(#3\right)} % \bias (auto-sizing)
}

% MSE (optional subscript, blank by default)
\NewDocumentCommand{\MSE}{ s O{} m }{
    \IfBooleanTF{#1}
        {\operatorname{MSE}_{#2}(#3)}              % \MSE*{} (no resizing)
        {\operatorname{MSE}_{#2}\!\left(#3\right)} % \MSE (auto-sizing)
}

% "full" derivative 
\NewDocumentCommand{\dx}{ o }{
    \mathrm{d}
    \IfValueT{#1}{
        \IfBlankF{#1}{#1}
    }
}

% "full" derivative (fractional)
\NewDocumentCommand{\ddx}{ O{} m m }{
    \IfValueTF{#1}
        {\frac{\dx^{#1}#2}{\dx{#3^{#1}}}}
        {\frac{\dx{#2}}{\dx{#3}}}
}

% partial derivative
\NewDocumentCommand{\px}{ o }{
    \partial
    \IfValueT{#1}{
        \IfBlankF{#1}{#1}
    }
}

% partial derivative (fractional)
\NewDocumentCommand{\pdx}{ O{} m m }{
    \IfValueTF{#1}
        {\frac{\px^{#1}#2}{\px{#3^{#1}}}}
        {\frac{\px{#2}}{\px{#3}}}
}

% inner product (optional base, blank by default)
\NewDocumentCommand{\innerprod}{ s O{} m m }{
    \IfBooleanTF{#1}
        {\langle{#3},{#4}\rangle_{#2}}            % \innerprod*{} (no resizing)       
        {\left\langle{#3},{#4}\right\rangle_{#2}} % \innerprod{}  (auto-sizing)
}

% norm (optional base, blank by default)
\NewDocumentCommand{\norm}{ s O{} m }{
    \IfBooleanTF{#1}
        {\lVert{#3}\rVert_{#2}}            % \norm*{} (no resizing)
        {\left\lVert{#3}\right\rVert_{#2}} % \norm{}  (auto-sizing)
}

% absolute value
\NewDocumentCommand{\abs}{ s m }{
    \IfBooleanTF{#1}
        {\lvert{#2}\rvert}            % \abs*{} (no resizing)
        {\left\lvert{#2}\right\rvert} % \abs{}  (auto-sizing)
}
%%% packages and libraries %%%
\usepackage{mathtools}
\usepackage{amsmath, amsthm, amssymb, amsfonts}
\usepackage{bm}

%%% blackboard alphabet (uppercase) %%%
\newcommand{\bbA}{\mathbb{A}}
\newcommand{\bbB}{\mathbb{B}}
\newcommand{\bbC}{\mathbb{C}}
\newcommand{\bbD}{\mathbb{D}}
\newcommand{\bbE}{\mathbb{E}}
\newcommand{\bbF}{\mathbb{F}}
\newcommand{\bbG}{\mathbb{G}}
\newcommand{\bbH}{\mathbb{H}}
\newcommand{\bbI}{\mathbb{I}}
\newcommand{\bbJ}{\mathbb{J}}
\newcommand{\bbK}{\mathbb{K}}
\newcommand{\bbL}{\mathbb{L}}
\newcommand{\bbM}{\mathbb{M}}
\newcommand{\bbN}{\mathbb{N}}
\newcommand{\bbO}{\mathbb{O}}
\newcommand{\bbP}{\mathbb{P}}
\newcommand{\bbQ}{\mathbb{Q}}
\newcommand{\bbR}{\mathbb{R}}
\newcommand{\bbS}{\mathbb{S}}
\newcommand{\bbT}{\mathbb{T}}
\newcommand{\bbU}{\mathbb{U}}
\newcommand{\bbV}{\mathbb{V}}
\newcommand{\bbW}{\mathbb{W}}
\newcommand{\bbX}{\mathbb{X}}
\newcommand{\bbY}{\mathbb{Y}}
\newcommand{\bbZ}{\mathbb{Z}}

%%% calligraphic alphabet (uppercase) %%%
\newcommand{\calA}{\mathcal{A}}
\newcommand{\calB}{\mathcal{B}}
\newcommand{\calC}{\mathcal{C}}
\newcommand{\calD}{\mathcal{D}}
\newcommand{\calE}{\mathcal{E}}
\newcommand{\calF}{\mathcal{F}}
\newcommand{\calG}{\mathcal{G}}
\newcommand{\calH}{\mathcal{H}}
\newcommand{\calI}{\mathcal{I}}
\newcommand{\calJ}{\mathcal{J}}
\newcommand{\calK}{\mathcal{K}}
\newcommand{\calL}{\mathcal{L}}
\newcommand{\calM}{\mathcal{M}}
\newcommand{\calN}{\mathcal{N}}
\newcommand{\calO}{\mathcal{O}}
\newcommand{\calP}{\mathcal{P}}
\newcommand{\calQ}{\mathcal{Q}}
\newcommand{\calR}{\mathcal{R}}
\newcommand{\calS}{\mathcal{S}}
\newcommand{\calT}{\mathcal{T}}
\newcommand{\calU}{\mathcal{U}}
\newcommand{\calV}{\mathcal{V}}
\newcommand{\calW}{\mathcal{W}}
\newcommand{\calX}{\mathcal{X}}
\newcommand{\calY}{\mathcal{Y}}
\newcommand{\calZ}{\mathcal{Z}}

%%% standard alphabet (uppercase) %%%
\newcommand{\rmA}{\mathrm{A}}
\newcommand{\rmB}{\mathrm{B}}
\newcommand{\rmC}{\mathrm{C}}
\newcommand{\rmD}{\mathrm{D}}
\newcommand{\rmE}{\mathrm{E}}
\newcommand{\rmF}{\mathrm{F}}
\newcommand{\rmG}{\mathrm{G}}
\newcommand{\rmH}{\mathrm{H}}
\newcommand{\rmI}{\mathrm{I}}
\newcommand{\rmJ}{\mathrm{J}}
\newcommand{\rmK}{\mathrm{K}}
\newcommand{\rmL}{\mathrm{L}}
\newcommand{\rmM}{\mathrm{M}}
\newcommand{\rmN}{\mathrm{N}}
\newcommand{\rmO}{\mathrm{O}}
\newcommand{\rmP}{\mathrm{P}}
\newcommand{\rmQ}{\mathrm{Q}}
\newcommand{\rmR}{\mathrm{R}}
\newcommand{\rmS}{\mathrm{S}}
\newcommand{\rmT}{\mathrm{T}}
\newcommand{\rmU}{\mathrm{U}}
\newcommand{\rmV}{\mathrm{V}}
\newcommand{\rmW}{\mathrm{W}}
\newcommand{\rmX}{\mathrm{X}}
\newcommand{\rmY}{\mathrm{Y}}
\newcommand{\rmZ}{\mathrm{Z}}

%%% standard alphabet (lowercase) %%%
\newcommand{\rma}{\mathrm{a}}
\newcommand{\rmb}{\mathrm{b}}
\newcommand{\rmc}{\mathrm{c}}
\newcommand{\rmd}{\mathrm{d}}
\newcommand{\rme}{\mathrm{e}}
\newcommand{\rmf}{\mathrm{f}}
\newcommand{\rmg}{\mathrm{g}}
\newcommand{\rmh}{\mathrm{h}}
\newcommand{\rmi}{\mathrm{i}}
\newcommand{\rmj}{\mathrm{j}}
\newcommand{\rmk}{\mathrm{k}}
\newcommand{\rml}{\mathrm{l}}
\newcommand{\rmm}{\mathrm{m}}
\newcommand{\rmn}{\mathrm{n}}
\newcommand{\rmo}{\mathrm{o}}
\newcommand{\rmp}{\mathrm{p}}
\newcommand{\rmq}{\mathrm{q}}
\newcommand{\rmr}{\mathrm{r}}
\newcommand{\rms}{\mathrm{s}}
\newcommand{\rmt}{\mathrm{t}}
\newcommand{\rmu}{\mathrm{u}}
\newcommand{\rmv}{\mathrm{v}}
\newcommand{\rmw}{\mathrm{w}}
\newcommand{\rmx}{\mathrm{x}}
\newcommand{\rmy}{\mathrm{y}}
\newcommand{\rmz}{\mathrm{z}}

%%% boldface math alphabet (uppercase) %%%
\newcommand{\bfA}{\mathbf{A}}
\newcommand{\bfB}{\mathbf{B}}
\newcommand{\bfC}{\mathbf{C}}
\newcommand{\bfD}{\mathbf{D}}
\newcommand{\bfE}{\mathbf{E}}
\newcommand{\bfF}{\mathbf{F}}
\newcommand{\bfG}{\mathbf{G}}
\newcommand{\bfH}{\mathbf{H}}
\newcommand{\bfI}{\mathbf{I}}
\newcommand{\bfJ}{\mathbf{J}}
\newcommand{\bfK}{\mathbf{K}}
\newcommand{\bfL}{\mathbf{L}}
\newcommand{\bfM}{\mathbf{M}}
\newcommand{\bfN}{\mathbf{N}}
\newcommand{\bfO}{\mathbf{O}}
\newcommand{\bfP}{\mathbf{P}}
\newcommand{\bfQ}{\mathbf{Q}}
\newcommand{\bfR}{\mathbf{R}}
\newcommand{\bfS}{\mathbf{S}}
\newcommand{\bfT}{\mathbf{T}}
\newcommand{\bfU}{\mathbf{U}}
\newcommand{\bfV}{\mathbf{V}}
\newcommand{\bfW}{\mathbf{W}}
\newcommand{\bfX}{\mathbf{X}}
\newcommand{\bfY}{\mathbf{Y}}
\newcommand{\bfZ}{\mathbf{Z}}

%%% boldface math alphabet (lowercase) %%%
\newcommand{\bfa}{\mathbf{a}}
\newcommand{\bfb}{\mathbf{b}}
\newcommand{\bfc}{\mathbf{c}}
\newcommand{\bfd}{\mathbf{d}}
\newcommand{\bfe}{\mathbf{e}}
\newcommand{\bff}{\mathbf{f}}
\newcommand{\bfg}{\mathbf{g}}
\newcommand{\bfh}{\mathbf{h}}
\newcommand{\bfi}{\mathbf{i}}
\newcommand{\bfj}{\mathbf{j}}
\newcommand{\bfk}{\mathbf{k}}
\newcommand{\bfl}{\mathbf{l}}
\newcommand{\bfm}{\mathbf{m}}
\newcommand{\bfn}{\mathbf{n}}
\newcommand{\bfo}{\mathbf{o}}
\newcommand{\bfp}{\mathbf{p}}
\newcommand{\bfq}{\mathbf{q}}
\newcommand{\bfr}{\mathbf{r}}
\newcommand{\bfs}{\mathbf{s}}
\newcommand{\bft}{\mathbf{t}}
\newcommand{\bfu}{\mathbf{u}}
\newcommand{\bfv}{\mathbf{v}}
\newcommand{\bfw}{\mathbf{w}}
\newcommand{\bfx}{\mathbf{x}}
\newcommand{\bfy}{\mathbf{y}}
\newcommand{\bfz}{\mathbf{z}}

%%% bold italic alphabet (uppercase) %%%
\newcommand{\bmA}{\bm{A}}
\newcommand{\bmB}{\bm{B}}
\newcommand{\bmC}{\bm{C}}
\newcommand{\bmD}{\bm{D}}
\newcommand{\bmE}{\bm{E}}
\newcommand{\bmF}{\bm{F}}
\newcommand{\bmG}{\bm{G}}
\newcommand{\bmH}{\bm{H}}
\newcommand{\bmI}{\bm{I}}
\newcommand{\bmJ}{\bm{J}}
\newcommand{\bmK}{\bm{K}}
\newcommand{\bmL}{\bm{L}}
\newcommand{\bmM}{\bm{M}}
\newcommand{\bmN}{\bm{N}}
\newcommand{\bmO}{\bm{O}}
\newcommand{\bmP}{\bm{P}}
\newcommand{\bmQ}{\bm{Q}}
\newcommand{\bmR}{\bm{R}}
\newcommand{\bmS}{\bm{S}}
\newcommand{\bmT}{\bm{T}}
\newcommand{\bmU}{\bm{U}}
\newcommand{\bmV}{\bm{V}}
\newcommand{\bmW}{\bm{W}}
\newcommand{\bmX}{\bm{X}}
\newcommand{\bmY}{\bm{Y}}
\newcommand{\bmZ}{\bm{Z}}

%%% bold italic alphabet (lowercase) %%%
\newcommand{\bma}{\bm{a}}
\newcommand{\bmb}{\bm{b}}
\newcommand{\bmc}{\bm{c}}
\newcommand{\bmd}{\bm{d}}
\newcommand{\bme}{\bm{e}}
\newcommand{\bmf}{\bm{f}}
\newcommand{\bmg}{\bm{g}}
\newcommand{\bmh}{\bm{h}}
\newcommand{\bmi}{\bm{i}}
\newcommand{\bmj}{\bm{j}}
\newcommand{\bmk}{\bm{k}}
\newcommand{\bml}{\bm{l}}
\newcommand{\bmm}{\bm{m}}
\newcommand{\bmn}{\bm{n}}
\newcommand{\bmo}{\bm{o}}
\newcommand{\bmp}{\bm{p}}
\newcommand{\bmq}{\bm{q}}
\newcommand{\bmr}{\bm{r}}
\newcommand{\bms}{\bm{s}}
\newcommand{\bmt}{\bm{t}}
\newcommand{\bmu}{\bm{u}}
\newcommand{\bmv}{\bm{v}}
\newcommand{\bmw}{\bm{w}}
\newcommand{\bmx}{\bm{x}}
\newcommand{\bmy}{\bm{y}}
\newcommand{\bmz}{\bm{z}}

%%% sans-serif alphabet (uppercase) %%%
\newcommand{\sfA}{\mathsf{A}}
\newcommand{\sfB}{\mathsf{B}}
\newcommand{\sfC}{\mathsf{C}}
\newcommand{\sfD}{\mathsf{D}}
\newcommand{\sfE}{\mathsf{E}}
\newcommand{\sfF}{\mathsf{F}}
\newcommand{\sfG}{\mathsf{G}}
\newcommand{\sfH}{\mathsf{H}}
\newcommand{\sfI}{\mathsf{I}}
\newcommand{\sfJ}{\mathsf{J}}
\newcommand{\sfK}{\mathsf{K}}
\newcommand{\sfL}{\mathsf{L}}
\newcommand{\sfM}{\mathsf{M}}
\newcommand{\sfN}{\mathsf{N}}
\newcommand{\sfO}{\mathsf{O}}
\newcommand{\sfP}{\mathsf{P}}
\newcommand{\sfQ}{\mathsf{Q}}
\newcommand{\sfR}{\mathsf{R}}
\newcommand{\sfS}{\mathsf{S}}
\newcommand{\sfT}{\mathsf{T}}
\newcommand{\sfU}{\mathsf{U}}
\newcommand{\sfV}{\mathsf{V}}
\newcommand{\sfW}{\mathsf{W}}
\newcommand{\sfX}{\mathsf{X}}
\newcommand{\sfY}{\mathsf{Y}}
\newcommand{\sfZ}{\mathsf{Z}}

%%% sans-serif alphabet (lowercase) %%%
\newcommand{\sfa}{\mathsf{a}}
\newcommand{\sfb}{\mathsf{b}}
\newcommand{\sfc}{\mathsf{c}}
\newcommand{\sfd}{\mathsf{d}}
\newcommand{\sfe}{\mathsf{e}}
\newcommand{\sff}{\mathsf{f}}
\newcommand{\sfg}{\mathsf{g}}
\newcommand{\sfh}{\mathsf{h}}
\newcommand{\sfi}{\mathsf{i}}
\newcommand{\sfj}{\mathsf{j}}
\newcommand{\sfk}{\mathsf{k}}
\newcommand{\sfl}{\mathsf{l}}
\newcommand{\sfm}{\mathsf{m}}
\newcommand{\sfn}{\mathsf{n}}
\newcommand{\sfo}{\mathsf{o}}
\newcommand{\sfp}{\mathsf{p}}
\newcommand{\sfq}{\mathsf{q}}
\newcommand{\sfr}{\mathsf{r}}
\newcommand{\sfs}{\mathsf{s}}
\newcommand{\sft}{\mathsf{t}}
\newcommand{\sfu}{\mathsf{u}}
\newcommand{\sfv}{\mathsf{v}}
\newcommand{\sfw}{\mathsf{w}}
\newcommand{\sfx}{\mathsf{x}}
\newcommand{\sfy}{\mathsf{y}}
\newcommand{\sfz}{\mathsf{z}}

%%% bold italic sans-serif alphabet (uppercase) %%%
% \newcommand{\bsfA}{\mathbfsf{A}}
% \newcommand{\bsfB}{\mathbfsf{B}}
% \newcommand{\bsfC}{\mathbfsf{C}}
% \newcommand{\bsfD}{\mathbfsf{D}}
% \newcommand{\bsfE}{\mathbfsf{E}}
% \newcommand{\bsfF}{\mathbfsf{F}}
% \newcommand{\bsfG}{\mathbfsf{G}}
% \newcommand{\bsfH}{\mathbfsf{H}}
% \newcommand{\bsfI}{\mathbfsf{I}}
% \newcommand{\bsfJ}{\mathbfsf{J}}
% \newcommand{\bsfK}{\mathbfsf{K}}
% \newcommand{\bsfL}{\mathbfsf{L}}
% \newcommand{\bsfM}{\mathbfsf{M}}
% \newcommand{\bsfN}{\mathbfsf{N}}
% \newcommand{\bsfO}{\mathbfsf{O}}
% \newcommand{\bsfP}{\mathbfsf{P}}
% \newcommand{\bsfQ}{\mathbfsf{Q}}
% \newcommand{\bsfR}{\mathbfsf{R}}
% \newcommand{\bsfS}{\mathbfsf{S}}
% \newcommand{\bsfT}{\mathbfsf{T}}
% \newcommand{\bsfU}{\mathbfsf{U}}
% \newcommand{\bsfV}{\mathbfsf{V}}
% \newcommand{\bsfW}{\mathbfsf{W}}
% \newcommand{\bsfX}{\mathbfsf{X}}
% \newcommand{\bsfY}{\mathbfsf{Y}}
% \newcommand{\bsfZ}{\mathbfsf{Z}}

%%% bold italic sans-serif alphabet (lowercase) %%%
% \newcommand{\bsfa}{\mathbfsf{a}}
% \newcommand{\bsfb}{\mathbfsf{b}}
% \newcommand{\bsfc}{\mathbfsf{c}}
% \newcommand{\bsfd}{\mathbfsf{d}}
% \newcommand{\bsfe}{\mathbfsf{e}}
% \newcommand{\bsff}{\mathbfsf{f}}
% \newcommand{\bsfg}{\mathbfsf{g}}
% \newcommand{\bsfh}{\mathbfsf{h}}
% \newcommand{\bsfi}{\mathbfsf{i}}
% \newcommand{\bsfj}{\mathbfsf{j}}
% \newcommand{\bsfk}{\mathbfsf{k}}
% \newcommand{\bsfl}{\mathbfsf{l}}
% \newcommand{\bsfm}{\mathbfsf{m}}
% \newcommand{\bsfn}{\mathbfsf{n}}
% \newcommand{\bsfo}{\mathbfsf{o}}
% \newcommand{\bsfp}{\mathbfsf{p}}
% \newcommand{\bsfq}{\mathbfsf{q}}
% \newcommand{\bsfr}{\mathbfsf{r}}
% \newcommand{\bsfs}{\mathbfsf{s}}
% \newcommand{\bsft}{\mathbfsf{t}}
% \newcommand{\bsfu}{\mathbfsf{u}}
% \newcommand{\bsfv}{\mathbfsf{v}}
% \newcommand{\bsfw}{\mathbfsf{w}}
% \newcommand{\bsfx}{\mathbfsf{x}}
% \newcommand{\bsfy}{\mathbfsf{y}}
% \newcommand{\bsfz}{\mathbfsf{z}}

%%% miscellaneous useful symbol commands %%%
\newcommand{\bfbeta}{\bm{\beta}}
\newcommand{\bfepsilon}{\bm{\epsilon}}
\newcommand{\bfmu}{\bm{\mu}}
\newcommand{\bftheta}{\bm{\theta}}
\newcommand{\bfSigma}{\bm{\Sigma}}

\newcommand{\sfTsmall}{\scriptstyle{\mathsf{T}}}
\newcommand{\sfHsmall}{\scriptstyle{\mathsf{H}}}

\newtheorem{theorem}{Theorem}[section]
\newtheorem{definition}{Definition}[section]
\newtheorem{assumption}{Assumption}[section]
\numberwithin{equation}{section}
% \newtheorem{theorem}{Theorem}
% \newtheorem{definition}{Definition}
% \newtheorem{assumption}{Assumption}

\usepackage{hyperref}
\usepackage[nameinlink]{cleveref}
\crefname{assumption}{assumption}{assumptions}
\Crefname{assumption}{Assumption}{Assumptions}

\usepackage{url}
\usepackage[dvipsnames]{xcolor}

\definecolor{linkRed}{HTML}{ED1B23}
\definecolor{citeBlue}{HTML}{006EB8}
\definecolor{urlBlue}{HTML}{2B2BB0}
\hypersetup{
    colorlinks = true,
    pdfborder = {0 0 0},
    citebordercolor = {1 1 1},
    linkcolor = linkRed,
    citecolor = citeBlue,
    urlcolor = urlBlue
}

\usepackage{tocloft}
\tocloftpagestyle{empty}

\newcommand{\PM}[1]{\textrm{PM}_{#1}}
\newcommand{\calGP}{\mathcal{GP}}
\DeclareMathOperator{\vect}{vec}

\title{\textbf{Computationally Efficient Predictions for Multi-Output Spatiotemporal Dynamical Systems}} % put your fancy title here

\author{
    Joey Hotz \\ \vspace{-0.3cm}\normalsize\texttt{joey.hotz@mail.utoronto.ca}\\ % put your beautiful name and email here
    \large
    \vspace{0.5cm}
    Supervisor: Dr. Meredith Franklin \\ % names of your almighty supervisors

    Examiners: Dr. Ricardo Baptista \& Dr. Joshua Speagle

    \vspace{0.5cm}
    Department of Statistical Sciences \\
    University of Toronto
}

% \author{
% Student Name \\ \texttt{emailaddress@mail.utoronto.ca} \\
% \and
% Supervisor(s): Supervisor Name(s) \\ [.5cm]Examiners: Examiner Names
% }
% \author{Student Name\thanks{Email: \texttt{emailaddress@mail.utoronto.ca}. Supervisor(s): Supervisor Name(s). Examiners: Examiner Names.}}
% \author{supervised by: Supervisor A}


% \affil{Department of Statistical Sciences\\University of Toronto}
% \affil[1]{Department of Statistical Sciences, University of Toronto, Canada}
% \affil[2]{Any additional affiliation}
% \affil[3]{Any additional affiliation}

\date{\today}

\onehalfspacing

\begin{document}

\maketitle
\thispagestyle{empty}

\begin{abstract}
    Over the past decade, several groups of machine learning researchers have demonstrated promising computational algorithms and mathematical results which utilize Gaussian process (GP) regression and Kalman filtering (KF) to derive estimates and predictions for the future states of spatiotemporal dynamical systems in a more computationally-efficient manner.
    
    This paper builds on previous results in this research area by adapting and extending the results of previous work on univariate outputs towards dynamical systems with vector-valued outputs. In this paper, we introduce the mathematical background necessary to understand Gaussian processes and Kalman filtering, and cover prior results in these areas. We then formally describe the multi-output spatiotemporal dynamics problem, and depict the assumptions required for similar computational speedups to exist in the multi-output spatiotemporal dynamics setting. Lastly, we apply these algorithms to both synthetic and real-world datasets to evaluate their performance.
\end{abstract}

% \bigskip
\clearpage

\tableofcontents

\newpage

\setcounter{page}{1}

\section{Introduction}
\label{sec:intro}

Gaussian processes (GPs) are an `industry standard' procedure for modelling the values of unknown functions, especially in scenarios where there are noisy or otherwise unreliable measurements and/or correlations in the data which must be accounted for \citep{GPML-Book-2006}. 

Gaussian process regression, also commonly referred to in the literature as \emph{kriging} after the South African geostatistician Danie G. Krige \citep{Cressie-Origins-1990}, is a statistical modelling procedure wherein a finite collection of observations are used for learning the form of an unknown real-valued function. Gaussian processes are particularly useful when there are correlations among observations, as these models can leverage information from correlated data to reduce the variance of outputs in the estimated function. Furthermore, the equations for fitting a GP to a dataset have computable closed-form solutions. Due to these convenient properties, Gaussian processes are frequently used in a variety of practical statistical applications involving modelling spatiotemporal data, particularly when we aim to learn certain functional properties of the underlying data.

Although Gaussian processes are mathematically tractable and convenient to use, there is one \emph{major} drawback to these models: exact inference with Gaussian processes has a cubically-scaling computational overhead. Fitting a Gaussian process onto a dataset which contains $M$ observations is a procedure requiring $\calO(M^{3})$ floating point operations \citep{GPML-Book-2006}, making these regression problems infeasible for many large-scale datasets \citep{Guhaniyogi-Meta-Kriging-2018,Pfoertner-Comp-Aware-2025,Salari-GP-Turbines-2025}. For instance, fitting a Gaussian process on a `moderately-sized' dataset of 1000 observations would involve $\calO(10^{9})$ floating point operations. In applied scenarios involving spatiotemporal modelling where there are repeated measurements, the complexity of Gaussian process regression will scale with the cube of the number of locations \emph{and} measurements per location, which may quickly extend beyond the scope of feasible calculation.

Furthermore, statistical modelling with Gaussian processes requires choosing a spatiotemporal covariance model which yields a positive semi-definite matrix on \emph{any} set of inputs \citep{GPML-Book-2006}. The class of valid covariance functions for a Gaussian process are mathematically convenient to work with in such regression contexts, but Gaussian process regression does impose further restrictions on how we choose to represent the covariance among observations in a spatiotemporal data setting. 

Due to these computational difficulties, many statisticians and machine learning researchers have sought to develop algorithms for fitting a Gaussian process with reduced computational overhead. Prominent well-studied and broadly applicable methods for accomplishing this task include approximations using induced latent variables \citep{Quinonero-Sparse-GP-2005}, sparse variational Gaussian processes \citep{Titsias-SVGP-2009}, and Vecchia approximations \citep{Katzfuss-Vecchia-2021}. However, these methods all require some degree of approximation, trading the exactness of a solution for increased algorithmic efficiency. 

In their 2020 paper \emph{Efficient spatio-temporal Gaussian regression via Kalman filtering}, \citeauthor{Todescato-Efficient-2020} introduce a new algorithm for efficiently deriving estimates of the future states in a spatio-temporal dynamical system under certain assumptions for the covariance structure between observations. The algorithm proposed by the authors utilizes Gaussian process regression in tandem with the Kalman filter, providing the ability to account for both the spatial and temporal aspects which guide the `evolution' of the dynamical system while simultaneously reducing the overall computational load. 

This paper provides a detailed overview of the results in \citet{Todescato-Efficient-2020} and subsequently builds upon the scaffolding in the authors' work by amending their proposed algorithm to operate with vector-valued outputs instead of being confined to scalar responses. This extension to multivariate response vectors provides many practical benefits, as there are many scientific research areas ranging from geostatistics and climatology to fields such as electrical engineering and signal processing \citep{Sarkka-Spatial-Kalman-2012,Huang-Nonstationary-2025} which utilize spatiotemporal models to account for the relationships between correlated response variables.

\section{Background}
\label{sec:background}

In order to properly describe the algorithm devised by \citet{Todescato-Efficient-2020}, we must first separately discuss Gaussian processes and Kalman filtering.

\subsection{Gaussian Processes}
\label{subsec:intro-GP}

A Gaussian process (GP) is a real-valued stochastic process indexed over some $P$-dimensional \emph{state space} $\calX \subseteq \bbR^{P}$ which satisfies the following key property:

\begin{definition}[Gaussian Process]
\label{def:GP}
    The real-valued function $f: \calX \to \bbR$ follows a Gaussian process on $\calX$ if for every finite collection ($N$-tuple) of states $\{\bfx_{1}, \cdots, \bfx_{N}\} \subset \calX$, the joint distribution of the random vector $\big(f(\bfx_{1}), \cdots, f(\bfx_{N})\big)^{\sfTsmall}$ is a multivariate Gaussian distribution on $\bbR^{N}$. 
\end{definition}

Furthermore, every Gaussian process on the state space $\calX$ is \emph{uniquely} specified by two real-valued functions: a mean function $m: \calX \to \bbR$ and a \emph{positive-definite} covariance (kernel) function $k: \calX \times \calX \to \bbR^{\ge 0}$, which respectively describe the mean and covariance of the underlying process \citep{GPML-Book-2006}. Any function $f: \calX \to \bbR$ which is a realization of a Gaussian process with mean $m(\cdot)$ and kernel $k(\cdot, \cdot)$ is endowed with the following structure:
\begin{subequations}
    \label{eq:GP-prior}
    \begin{align}
        \textbf{Mean:} \quad & m(\bfx) = \expec{f(\bfx)}\; \textrm{for all}\; \bfx \in \calX\label{eq:GP-prior-mean}\\% [6pt]
        \textbf{Covariance:} \quad & k(\bfx, \bfx') = \covar{f(\bfx)}{f(\bfx')}\; \textrm{for all}\;(\bfx, \bfx') \in \calX \times \calX\label{eq:GP-prior-kernel}
    \end{align}
\end{subequations}
We denote a function $f: \calX \to \bbR$ which is drawn from a Gaussian process with mean $m(\cdot)$ and kernel $k(\cdot, \cdot)$ as $f \sim \calGP(m, k)$. Note that the kernel function $k(\cdot, \cdot)$ in \Cref{eq:GP-prior-kernel} must be symmetric in its arguments, as this function describes the covariance between the Gaussian process outputs at these two locations. In addition to describing the covariance between function observations, the kernel function $k(\cdot, \cdot)$ serves another complementary purpose, as sample paths have continuous mean-square derivatives up to the order of the kernel $k(\cdot, x)$ \citep{GPML-Book-2006}.

Suppose that $f \sim \calGP(m, k)$ is a realization of some Gaussian process on $\calX \subseteq \bbR^{P}$, and let $\{\bfx_{1}, \dots, \bfx_{N}\} \subset \calX$ be a set of queried locations. For ease of notation, we define the data matrix $\bfX := \begin{bsmallmatrix}
    \bfx_{1}^{\scriptscriptstyle{\sfT}} &
    \cdots &
    \bfx_{N}^{\scriptscriptstyle{\sfT}}
\end{bsmallmatrix}^{\sfTsmall} \in \bbR^{N \times P}$, and define the corresponding vectors $f(\bfX) := \begin{bsmallmatrix}
    f(\bfx_{1}^{\scriptscriptstyle{\sfT}}) &
    \cdots &
    f(\bfx_{N}^{\scriptscriptstyle{\sfT}})
\end{bsmallmatrix}^{\sfTsmall} \in \bbR^{N}$ and $m(\bfX) := \begin{bsmallmatrix}
    m(\bfx_{1}^{\scriptscriptstyle{\sfT}}) &
    \cdots &
    m(\bfx_{N}^{\scriptscriptstyle{\sfT}})
\end{bsmallmatrix}^{\sfTsmall} \in \bbR^{N}$. Lastly, let $k(\bfX, \bfX) \in \bbR^{N \times N}$ represent the \emph{Gram matrix} for the observed locations which is induced by the kernel $k(\cdot, \cdot)$:
\begin{equation}
\label{eq:GP-gram-matrix}
    k(\bfX, \bfX) := \begin{bmatrix}
        k(\bfx_{1}, \bfx_{1}) & k(\bfx_{1}, \bfx_{2}) & \cdots & k(\bfx_{1}, \bfx_{N})\\
        k(\bfx_{2}, \bfx_{1}) & k(\bfx_{2}, \bfx_{2}) & \cdots & k(\bfx_{2}, \bfx_{N})\\
        \vdots & \vdots & \ddots & \vdots\\
        k(\bfx_{N}, \bfx_{1}) & k(\bfx_{N}, \bfx_{2}) & \cdots & k(\bfx_{N}, \bfx_{N})
    \end{bmatrix}
\end{equation}
One of the key properties which allows Gaussian processes to be so widely-applicable is that under the conditions described above, the vector $f(\bfX)$ follows a multivariate Gaussian distribution:
\begin{equation}
\label{eq:GP-fun-distribution}
    f(\bfX) \sim \calN\Big(m(\bfX), k(\bfX, \bfX)\Big)
\end{equation}
Furthermore, consider the same setting as above, and suppose that we have some (potentially) noisy measured observations given by $y_{i} = f(\bfx_{i}) + \epsilon_{i}$, for independent normally-distributed measurement errors $\epsilon_{1}, \dots, \epsilon_{N} \simiid \calN(0, \sigma^2)$. Equivalently, we assume that $y_{i} \sim \calN\big(f(\bfx_{i}), \sigma^{2}\big)$ for each queried location $\bfx_{i}$ \citep{GPML-Book-2006}. Defining the response vector $\bfy := \begin{bsmallmatrix}
    y_{1}^{\phantom{{\scriptscriptstyle{\sfT}}}} &
    \cdots &
    y_{N}^{\phantom{{\scriptscriptstyle{\sfT}}}}
\end{bsmallmatrix}^{\sfTsmall} \in \bbR^{N}$ yields the following distributional result:
\begin{equation}
\label{eq:GP-obs-distribution}
    \bfy \sim \calN\Big(m(\bfX), k(\bfX, \bfX) + \sigma^{2}\bfI_{N}\Big)
\end{equation}
Gaussian processes are often incorporated in a Bayesian setting as a prior distribution which one can place on the space of real-valued functions on $\calX$ \citep{Titsias-SVGP-2009,Alvarez-Vector-Kernels-2012}. For example, for a \emph{black-box} function $f: \calX \to \bbR$, we may place a prior of the form $f \sim \calGP(m, k)$, for specified functions $m: \calX \to \bbR$ and $k: \calX \times \calX \to \bbR$ as described by \Cref{eq:GP-prior-mean,eq:GP-prior-kernel}. If observations of the black-box function $f(\cdot)$ are taken with normally-distributed measurement error (as described by \Cref{eq:GP-obs-distribution}), then the posterior predictive distribution for $f(\cdot)$ will be a Gaussian process whose mean and kernel are altered by conditioning on the observed data. In particular, for a data matrix $\bfX \in \bbR^{N \times P}$ with the response vector $\bfy \in \bbR^{N}$, the posterior predictive distribution for $f$ evaluated at any points $\bfx, \bfx' \in \calX$ is given by the following:
\begin{subequations}
    \label{eq:GP-posterior}
    \begin{align}
        \textbf{Mean:} \quad & m_{\textrm{post}}(\bfx) = m(\bfx) + k(\bfx, \bfX)\big[k(\bfX, \bfX) + \sigma^{2}\bfI_{N}\big]^{-1}\big(\bfy - m(\bfX)\big)\label{eq:GP-post-mean}\\% [6pt]
        \textbf{Kernel:} \quad & k_{\textrm{post}}(\bfx, \bfx') = k(\bfx, \bfx') - k(\bfx, \bfX)\big[k(\bfX, \bfX) + \sigma^{2}\bfI_{N}\big]^{-1}k(\bfX, \bfx')\label{eq:GP-post-kernel}
    \end{align}
\end{subequations}
In the equations above, $k(\bfx, \bfX)$ represents the $1$-by-$N$ Gram matrix with entries $k(\bfx, \bfx_{i})$ and similarly, $k(\bfX, \bfx')$ is the $N$-by-1 Gram matrix with entries $k(\bfx_{i}, \bfx')$. Setting $\bfx = \bfx'$ in \Cref{eq:GP-post-kernel} yields the posterior predictive variance for $f(\bfx)$. Furthermore, when the observations $y_{1}, \dots, y_{N}$ are assumed to be exact (i.e. $\sigma^2 = 0$), then \Cref{eq:GP-post-mean,eq:GP-post-kernel} will perfectly interpolate every observation. To be precise, if $\sigma^{2} = 0$, the posterior Gaussian process satisfies that $m_{\textrm{post}}(\bfx_{i}) = y_{i}$ and $k_{\textrm{post}}(\bfx_{i}, \bfx_{i}) = 0$, implying that each $f(\bfx_{i})$ is \emph{exactly} equal to the corresponding $y_{i}$. 

\subsection{Multi-Output Gaussian Processes}
\label{subsec:intro-MOGP}
Thus far, we have focused solely on the description of single-output Gaussian processes, whose realizations are described by functions of the form $f: \calX \to \bbR$. However, there is a natural extension from these univariate processes to \emph{multi-output} (also called \emph{multi-task}) Gaussian processes: a higher-dimensional form of a Gaussian process, whose outputs are vectors in $\bbR^{D}$ \citep{Alvarez-Vector-Kernels-2012,Bonilla-Multitask-GP-2007}. 

\begin{definition}[Multi-Output Gaussian Process]
\label{def:MOGP}
    The vector-valued function $\bff = (f_{1}, \dots, f_{D})^{\sfTsmall}: \calX \to \bbR^{D}$ follows a multi-output Gaussian process on $\calX$ if for every finite collection of states $\{\bfx_{1}, \dots, \bfx_{N}\} \subset \calX$, the joint distribution of the $ND$ random variables $\big(f_{1}(\bfx_{1}), \cdots, f_{D}(\bfx_{N})\big)^{\sfTsmall}$ is a multivariate Gaussian distribution on $\bbR^{ND}$.
\end{definition}

Multi-output GPs (MOGPs) are naturally applicable whenever there are several (potentially) correlated quantities which are measured over some common domain. These processes include both \emph{auto-correlations} within the same output and \emph{cross-correlations} between outputs, allowing for information about one output to reduce the uncertainty of other response variables in the model. Conceptually, a multi-output GP is analogous to simultaneously modelling $D$ single-output Gaussian processes, while also modelling the covariance between these processes.

Given states $\bfX = \begin{bsmallmatrix} \bfx_{1}^{\scriptscriptstyle{\sfT}} & \cdots & \bfx_{N}^{\scriptscriptstyle{\sfT}} \end{bsmallmatrix}^{\sfTsmall} \in \bbR^{N \times P}$, we define the matrix $\bff(\bfX) \in \bbR^{N \times D}$ by concatenating the vectorized columns, i.e. $\bff(\bfX) = \begin{bsmallmatrix}f_{1}(\bfX) & \cdots & f_{D}(\bfX)\end{bsmallmatrix}$, and we can analogously define the matrix $\bfm(\bfX) \in \bbR^{N \times D}$ as a concatenation of the respective $N$-dimensional mean vectors $m_{1}(\bfX), \dots, m_{D}(\bfX)$. In order to work with a multi-output GP, we must also construct a Gram matrix $\bfK(\bfX, \bfX) \in \bbR^{ND \times ND}$ which defines the covariance structure between both the queried locations and the output dimensions. In its most general form, the Gram matrix providing the covariance for a multi-output GP has the following structure:
\begin{equation}
\label{eq:MOGP-gram-matrix}
    \bfK(\bfX, \bfX) := \begin{bmatrix}
        \bfK(\bfx_{1}, \bfx_{1}) & \bfK(\bfx_{1}, \bfx_{2}) & \cdots & \bfK(\bfx_{1}, \bfx_{N})\\
        \bfK(\bfx_{2}, \bfx_{1}) & \bfK(\bfx_{2}, \bfx_{2}) & \cdots & \bfK(\bfx_{2}, \bfx_{N})\\
        \vdots & \vdots & \ddots & \vdots\\
        \bfK(\bfx_{N}, \bfx_{1}) & \bfK(\bfx_{N}, \bfx_{2}) & \cdots & \bfK(\bfx_{N}, \bfx_{N})
    \end{bmatrix}
\end{equation}
The covariance matrix in \Cref{eq:MOGP-gram-matrix} is overall quite similar to the Gram matrix in the univariate setting which is shown in \Cref{eq:GP-gram-matrix}. However, the multi-output Gram matrix consists of $N^{2}$ covariance matrix sub-blocks $\bfK(\bfx_{i}, \bfx_{j}) = \covar{\bff(\bfx_{i})}{\bff(\bfx_{j})}$ of size $D \times D$, each of which has entries of the form $[\bfK(\bfx_{i}, \bfx_{j})]_{(d,d')} = \covar{f_{d}(\bfx_{i})}{f_{d'}(\bfx_{j})}$ \citep{Alvarez-Vector-Kernels-2012}. Vectorizing the matrix $\bff(\bfX)$ yields the following distributional result:
\begin{equation}
\label{eq:MOGP-fun-distribution}
    \vect\!\big(\bff(\bfX)\big) \sim \calN\Big(\vect\!\big(\bfm(\bfX)\big),\, \bfK(\bfX, \bfX)\Big)
\end{equation}
Additionally, suppose that we have observations of the form $\bfy_{i} = \bff(\bfx_{i}) + \bfepsilon_{i}$ with independent measurement error vectors $\bfepsilon_{1}, \dots, \bfepsilon_{N} \simiid \calN(\mathbf{0}, \bfSigma_{\epsilon})$, for some positive-definite covariance matrix $\bfSigma_{\epsilon} \in \bbR^{D \times D}$. The joint distribution of these observed values is given by
\begin{equation}
\label{eq:MOGP-obs-distribution}
    \vect\!\big(\bff(\bfX)\big) \sim \calN\Big(\vect\!\big(\bfm(\bfX)\big),\, \bfK(\bfX, \bfX) + \bfI_{N} \otimes \bfSigma_{\epsilon}\Big)
\end{equation}
Again, the full distribution in \Cref{eq:MOGP-obs-distribution} closely echoes the univariate result from \Cref{eq:GP-obs-distribution}, with a slight adjustment to account for covariances between outputs.

From a purely theoretical standpoint, multi-output Gaussian processes are defined quite analogously to the single-output processes discussed in \Cref{subsec:intro-GP}. However, by solely introducing the theory and notation regarding MOGPs, we have obscured the major practical implications of these models; namely, the extreme cost for computational inference. In the (general) single-output Gaussian process regime, the cost of inference scales cubically with the size of the matrix $k(\bfX, \bfX)$, giving a computational cost of the order $\calO(N^{3})$. In the setting where the responses are vector-valued, statistical inference still requires the inversion of the Gram matrix $\bfK(\bfX, \bfX)$, whose cost balloons to the order of $\calO(N^{3}D^{3})$. Even when $D \ll N$, this is a significant increase in the computational requirement for an already-expensive operation. We will now turn our attention towards three specific cases of multi-output Gaussian processes, which give simpler forms for the Gram matrix $\bfK(\bfX, \bfX)$.

\paragraph{Fully Independent Model:} The simplest form of a multi-output Gaussian process is one in which each output dimension is assumed to be an independent univariate Gaussian process: $f_{d} \sim \calGP(m_{d}, k_{d})$. Under this assumption, the Gram matrix $\bfk(\bfX, \bfX)$ is a block-diagonal matrix with $D$ $N \times N$ sub-blocks. This model is computationally quite cheap, as it can be simply decoupled into $D$ separate single-output regression problems. However, this model makes the strong assumption that the targeted outputs are fully independent, and it shares no information across the output dimensions, which hampers its applicability in many practical contexts.

This model is computationally cheap, since it decouples into $D$ separate single-output problems, but it shares no information across outputs whatsoever. The models below reintroduce dependencies between the outputs by expressing them as the linear combinations of latent Gaussian processes.

\paragraph{Linear Model of Coregionalization (LMC):} The LMC model for a Gaussian process with a vector-valued output is predicated on the assumption that the output functions are built from a linear combination of some set of independent mean-zero univariate Gaussian processes \citep{Alvarez-Vector-Kernels-2012}. Suppose that the latent Gaussian processes $u_{1}, \dots, u_{R}$ are independent, with $\expec{u_{i}(\bfx)} = 0$ and $\covar{u_{r}(\bfx)}{u_{r'}(\bfx')} = \mathbf{1}_{\{r = r'\}} \times k_{r}(\bfx, \bfx')$, so that the covariance is nonzero if and only if $r = r'$. 

Following the description given by \citet{Alvarez-Vector-Kernels-2012}, we can sort these latent GPs based on their covariance function, giving us $Q$ groups of size $\{R_{q}\}_{q=1}^{Q}$. By sorting the latent processes in this manner, each group $q$ (for $q = 1, \dots, Q$) consists of $R_{q}$ independent univariate Gaussian processes of the form $u_{q}^{1}, \dots, u_{q}^{R_{q}} \simiid \calGP(0, K^{(q)})$. Note that even though the covariance structures are shared for latent functions in the same group, the respective Gaussian processes are still entirely independent. As such, $\covar*{u_{q}^{i}(\bfx)}{u_{q'}^{i'}(\bfx')} = \mathbf{1}_{\{q = q'\}} \times \mathbf{1}_{\{i = i'\}} \times K^{(q)}(\bfx, \bfx')$ under this regrouping of the latent functions.

In the LMC model, we assume that the desired functions $\{f_{d}(\cdot)\}_{d=1}^{D}$ for the multivariate process are constructed via a linear combination of the latent Gaussian processes $u_{1}, \dots, u_{R}$. Each function is therefore described by the following latent model:
\begin{align}
    f_{d}(\bfx) &= m_{d}(\bfx) + \sum_{r=1}^{R}\alpha_{d,r}u_{r}(\bfx) \label{eq:LMC-generative-ungrouped}
    \intertext{Grouping the latent functions by their kernel, this can be written in the form}
    f_{d}(\bfx) &= m_{d}(\bfx) + \sum_{q=1}^{Q}\sum_{i=1}^{R_{q}}\alpha_{d,q}^{i}u_{q}^{i}(\bfx) \label{eq:LMC-generative-grouped}
\end{align}
The functions $\{m_{d}(\cdot)\}_{d=1}^{D}$ in \Cref{eq:LMC-generative-ungrouped,eq:LMC-generative-grouped} ensure that each function in the observable multi-output process can have a non-zero mean. 
Using the latent form given in \Cref{eq:LMC-generative-grouped}, we can represent the covariance between two outputs by the following:
\begin{subequations}
    \label{eq:LMC-kernel}
    \begin{align}
        \covar{f_{d}(\bfx)}{f_{d'}(\bfx')} &= \sum_{q=1}^{Q}\sum_{q'=1}^{Q}\sum_{i=1}^{R_{q}}\sum_{i=1}^{R_{q'}}\alpha_{d,q}^{i}\alpha_{d',q'}^{i'}\covar{u_{q}^{i}(\bfx)}{u_{q'}^{i'}(\bfx)} \label{eq:LMC-kernel-long}
        \intertext{In the summation above, the covariance is non-zero if and only if $q = q'$ and $i = i'$.}
        &= \sum_{q=1}^{Q}\sum_{i=1}^{R_{q}}\alpha_{d,q}^{i}\alpha_{d',q}^{i}\covar{u_{q}^{i}(\bfx)}{u_{q}^{i}(\bfx)}\label{eq:LMC-kernel-shorter}\\
        &= \sum_{q=1}^{Q}\sum_{i=1}^{R_{q}}\alpha_{d,q}^{i}\alpha_{d',q}^{i}K^{(q)}(\bfx, \bfx') \label{eq:LMC-kernel-short}
    \end{align}
\end{subequations}
For each group $q = 1, \dots, Q$, and each pair of dimensions $(d, d')$, we define the value $b^{(q)}_{d, d'} := \sum_{i=1}^{R_{q}}\alpha_{d,q}^{i}\alpha_{d',q}^{i}$. We can then concatenate these values into a $D \times D$ \emph{coregionalization} matrix $\bfB^{(q)}$ such that the $(d,d')$ entry of $\bfB^{(q)} = b^{(q)}_{d,d'}$. Using these coregionalization matrices yields the kernel function for the linear model of coregionalization, which is given by 
\begin{equation}
    \bfK(\bfx, \bfx') = \sum_{q=1}^{Q}K^{(q)}(\bfx, \bfx')\bfB^{(q)}, \quad \textrm{with}\ \bfK(\bfx, \bfx')\in \bbR^{D \times D},
\end{equation}
and the overall Gram matrix for the linear model of coregionalization is given by
\begin{equation}
    \label{eq:LMC-gram}
    \bfK(\bfX, \bfX) = \sum_{q=1}^{Q}K^{(q)}(\bfX, \bfX) \otimes \bfB^{(q)}, \quad \textrm{with}\ \bfK(\bfX, \bfX) \in \bbR^{ND \times ND}
\end{equation}
The linear model of coregionalization is commonly used for describing the kernels of multi-output Gaussian processes in scientific modelling settings such as inference among sensor networks, geostatistical applications, and robot dynamics \citep{Alvarez-Vector-Kernels-2012, Liu-MOGP-Remarks-2018}.

\paragraph{Intrinsic Coregionalization Model (ICM):} The intrinsic coregionalization model (ICM) for a multi-output Gaussian process is a special case of the LMC kernel, where the number of distinct kernel functions ($Q$) is equal to one. The ICM kernel is a more restrictive simplification of the LMC kernel, but the Gram matrix will be significantly less expensive to decompose as a result of this simplifying assumption.

The ICM kernel can be derived in the same format as the LMC kernel, with all latent functions $u_{1}, \dots, u_{R} \simiid \calGP(0,k)$. The covariances generated by the ICM kernel are described by the following relations:
\begin{equation}
    \label{eq:ICM-covariance}
    \covar{f_{d}(\bfx)}{f_{d'}(\bfx')} = k(\bfx, \bfx') \times \bfB(d, d')
\end{equation}
The covariance given in \Cref{eq:ICM-covariance} corresponds to a kernel which is similar to that in \Cref{eq:LMC-gram}, with only one spatial kernel $k(\bfX, \bfX) \in \bbR^{N \times N}$ and one corresponding coregionalization matrix $\bfB \in \bbR^{D \times D}$ for the outputs. 
\begin{equation}
    \label{eq:ICM-gram}
    \bfK(\bfX, \bfX) = k(\bfX, \bfX) \otimes \bfB, \quad \textrm{with}\ \bfK(\bfX, \bfX) \in \bbR^{ND \times ND}
\end{equation}
The ICM kernel is an example of a \emph{separable} kernel, as the Gram matrix in \ref{eq:ICM-gram} can be written as the Kronecker product of one purely spatial Gram matrix and another purely output-related Gram matrix. Similarly, the individual terms of the form $k(\bfX, \bfX) \otimes \bfB^{(q)}$ in the summation in \ref{eq:LMC-gram} are separable kernels.

\paragraph{Semiparametric Latent Factor Model (SLFM):} Another special case of the LMC kernel is the \emph{semiparametric latent factor model} (SLFM) \citep{Bonilla-Multitask-GP-2007}, in which there is exactly one kernel corresponding to each group, i.e. $R_{q} = 1$ for each $q \in \{1, \dots, Q\}$ in the equation in \Cref{eq:LMC-generative-grouped}. Under this scenario, each matrix $\bfB^{(q)}$ is a rank-one matrix generated by the coefficients ${\alpha}_{d,q}$, which also allows for the resulting matrix decomposed in a less expensive manner than the traditional LMC kernel.

\paragraph{Other Constructions:} Many valid MOGP kernels do not satisfy the sum of separable kernels form which was depicted above, such as \emph{process convolution} approaches \citep{Boyle-Dependent-GP-2004} and other invariant forms of kernels which incorporate lags or spatial shifts \citep{Alvarez-Vector-Kernels-2012}. In this paper, we focus on multi-output kernels which can be represented via a sum of separable kernels, as the Kronecker product structure combines naturally with the separability requirements considered by \citet{Todescato-Efficient-2020}. This produces a convenient structure for the Gram matrices which we consider in \Cref{sec:problem}.

\subsection{Kalman Filtering}
\label{subsec:kalman}

The Kalman filter (KF), introduced and eponymously named for Rudolf E. K{\'a}lm{\'a}n \citeyearpar{Kalman-Filter-1960}, is a technique for recursively estimating the unobserved (latent) states of a linear dynamical system from a sequence of noisy and indirect measurements of the system. K{\'a}lm{\'a}n originally derived the filter as the means to solving an optimal linear estimation problem via orthogonal projections. However, the Kalman filter also lends itself to a statistical interpretation, as the filter is equivalent to the solution of a Bayesian inference problem for a \emph{linear Gaussian state-space model} whose posterior distribution is updated sequentially as new observations are recorded \citep{Bayesian-Filter-Book-2023}.

In this subsection of the paper, we will formally describe the Kalman filter, and the linear Gaussian state-space model which underpins the statistical viewpoint for the Kalman filter. These statistical properties of interest are central to the algorithm presented by \citet{Todescato-Efficient-2020}.

\begin{definition}[State Space Model]
\label{def:SSM}  
    Let $t_{0} < t_{1} < t_{2} < \cdots < t_{J}$ be an indexed sequence of time points. A state space model is a dynamical system describing the joint behaviour of two stochastic processes: a {latent state} process $\{\bfs_{j}\}_{j=0}^{J}$ for states $\bfs_{j} \in \bbR^{n}$ which evolves over time according to some Markovian dynamics, and an {observation} (or {measurement}) process $\{\bfy_{j}\}_{j=1}^{J}$, where each measurement $\bfy_{j} \in \bbR^{m}$ is a (potentially) noisy observation of some deterministic function of the current latent state $\bfs_{j}$.
\end{definition}

One particularly convenient form of a state space model is the \emph{linear Gaussian} state space model. In this setting, the dynamics which guide the evolution of the system over time and the observation model are both linear functions of the underlying states, and the random terms in the model follow multivariate Gaussian distributions.

\begin{definition}[Linear Gaussian State Space Model]
\label{def:LGSSM}
    A linear Gaussian state space model is a dynamical state space model whose latent state and observation model are described by the following equations:
    \begin{subequations}
        \label{eq:LGSSM-dynamics}
        \begin{align}
            \textnormal{\textbf{Initial state:}} \quad & \bfs(0) \sim \calN\big(\bfm(0), \bfSigma(0)\big) && \label{eq:LGSSM-init}\\% [6pt]
            \textnormal{\textbf{Dynamic model:}} \quad & \bfs(j+1) = \bfA(j)\bfs(j) + \bfw(j), \, && \bfw(j) \sim \calN\big(\mathbf{0}, \bfQ(j)\big) \label{eq:LGSSM-evolution}\\% [6pt]
            \textnormal{\textbf{Observation model:}} \quad & \bfy(j) = \bfC(j)\bfs(j) + \bfv(j), \, && \bfv(j) \sim \calN\big(\mathbf{0}, \bfR(j)\big) \label{eq:LGSSM-observe}
        \end{align}
    \end{subequations}
\end{definition}
In the definition above, $\{\bfs(j)\}_{j=0}^{J}$ are the latent states within the dynamical model. The evolution of the dynamical system over time is guided by the \emph{state transition matrix} $\bfA(j)$ and the random Gaussian noise $\bfw(j)$, and the observation process is influenced by the \emph{observation matrix} $\bfC(j)$ and observation noise $\bfv(j)$. The initial state $\bfs(0)$, process noise terms $\{\bfw(j)\}_{j=0}^{J-1}$ and observation noise terms $\{\bfv(j)\}_{j=1}^{J}$ are all assumed to be mutually independent. Furthermore, the parameters $\bfm(0), \bfSigma(0), \{\bfQ(j)\}_{j=0}^{J-1}$, and $\{\bfR(j)\}_{j=1}^{J}$, and the transition and observation matrices are assumed to have fixed, known values. 

A key property of the multivariate Gaussian distribution is that \emph{any} linear transformation -- including affine linear transformations -- of a Gaussian-distributed random vector will still follow a Gaussian distribution, as does the sum of any pair of \emph{independent} Gaussian-distributed vectors. Consequently, in the linear Gaussian state space model described by \Cref{def:LGSSM}, every latent state $\{\bfs(j)\}_{j=0}^{J}$ and measurement $\{\bfy(j)\}_{j=1}^{J}$ is a random vector drawn from a Gaussian distribution with mathematically tractable parameters \citep{Bayesian-Filter-Book-2023}.

In such cases, the standard method for recovering the parameters describing the distribution of each random vector is to utilize the Kalman filter \citep{Kalman-Filter-1960}, which iteratively computes the respective conditional mean and covariance for each latent state and observation. This is accomplished by iterating between two steps; the \emph{prediction} step and the \emph{update} step.


\begin{algorithm}
    \caption{The Kalman Filter \citep{Kalman-Filter-1960}}
    \label{alg:kalman-filter}
    
    \begin{algorithmic}[1]
        \Require
            \Statex \hspace{\algorithmicindent} $\widehat{\bfm}(0 \mid -1)$: expected value of the initial state $\bfs(0)$       
            \Statex \hspace{\algorithmicindent} $\bfSigma(0 \mid -1)$: variance of the initial latent state $\bfs(0)$
        \For{$j = 1$ to $J$}
            \Statex \texttt{\# Prediction Step:}
            \State $\widehat{\bfm}(j \mid j-1) = \bfA(j-1)\widehat{\bfm}(j-1 \mid j-1)$
            \State $\bfSigma(j \mid j-1) = \bfA(j-1)\bfSigma(j-1 \mid j-1)\bfA(j-1)^{\sfTsmall} + \bfQ(j-1)$

            \Statex \texttt{\# Update Step:}
            \State $\bfd(j) = \bfy(j) - \bfC(j)\widehat{\bfm}(j \mid j-1)$
            \State $\bfS(j) = \bfC(j)\bfSigma(j \mid j-1)\bfC(j)^{\sfTsmall} + \bfR(j)$
            \State $\bfK(j) = \bfSigma(j \mid j-1)\bfC(j)^{\sfTsmall}\bfS(j)^{-1}$ \Comment{Kalman gain}
            \State $\widehat{\bfm}(j \mid j) = \widehat{\bfm}(j \mid j-1) + \bfK(j)\bfd(j)$
            \State $\bfSigma(j \mid j) = \bfSigma(j \mid j-1) - \bfK(j)\bfS(j)\bfK(j)^{\sfTsmall}$
        \EndFor
    \end{algorithmic}
\end{algorithm}

The prediction and update steps in the Kalman filtering algorithm (\Cref{alg:kalman-filter}) are so named as they correspond to specific conditional distributions of $\bfs(j)$ and $\bfy(j)$ before and after conditioning on the latest observation. In particular, for the linear Gaussian state space model described by \Cref{def:LGSSM}, we have the following distributional relationships:
\begin{subequations}
    \label{eq:Kalman-conditionals}
    \begin{align}
        \bfs(j) \mid \bfy(1), \dots, \bfy(j-1) &\sim \calN\big(\widehat{\bfm}(j\mid j), \widehat{\bfSigma}(j)\big) \\% [6pt]
        \bfy(j) \mid \bfy(1), \dots, \bfy(j-1) &\sim \calN\big(\bfC(j)\widehat{\bfm}(j), \bfS(j)\big) \\% [6pt]
        \bfs(j) \mid \bfy(1), \dots, \bfy(j) &\sim \calN\big(\bfm(j), \bfSigma(j)\big)
    \end{align}
\end{subequations}
Each of the $J$ iterations of \Cref{alg:kalman-filter} requires both multiplying and inverting matrices.  In total, the computational complexity per iteration in the Kalman filtering loop is $\calO(m^{3} + n^{3})$ \citep{Kalman-Complexity-2019}, with each of these steps occurring $J$ times in the total run of the algorithm. Thus, a `naive' implementation of this algorithm would have a total computational complexity on the order of $\calO\big(J(m^{3} + n^{3})\big)$. Overall, the computational complexity of the Kalman filtering algorithm is \emph{linear} with respect to the number of iterations $J$, which is far quicker than the cubically-scaling $\calO(J^3)$ runtime required for Gaussian process regression on a set of $J$ observations. This algorithm also has a much lower storage requirement, as the updates are based strictly on the previous state, allowing for overwriting of earlier estimators.

\section{Problem Formulation}
\label{sec:problem}
Now that we have discussed the pre-requisite information regarding Gaussian process regression and Kalman filtering, we may begin to describe the setting considered by \citeauthor{Todescato-Efficient-2020}.

First, we will describe the initial univariate Gaussian process regression problem posed by the authors, and the algorithm which they propose as a means of more efficiently deriving their desired estimates. Following this, we will explain how the algorithm by \citeauthor{Todescato-Efficient-2020} can be extended for use with multi-output spatiotemporal dynamical systems, while maintaining similar computational benefits in comparison to related MOGP-based methods.

\subsection{Single-Output Dynamical System}
Suppose that we have some \emph{compact} state space $\calX \subseteq \bbR^{P}$ which represents a set of all possible locations of interest, such as the latitudes and longitudes of a bounded geographical region, or the coordinates of pixels in an image. Within the state space $\calX$, we have a \emph{finite} set of ``observable'' locations $\calI := \{x_{1}, \dots, x_{M}\} \subset \calX$.

We wish to perform inference on the values of some function $f: \calX \times \bbR^{+} \to \bbR$, which we assume to be a realization of some zero-mean Gaussian process $f \sim \calGP(0, K)$, for some kernel function $K$. Furthermore, at times $t_{0} < t_{1} < \dots$, we record noisy observations of the form 
\begin{equation}
    y_{i}(t_{j}) = f(x_{i}, t_{j}) + v_{i}(t_{j}), \quad v_{i}(t_{j}) \sim \calN(0, \sigma^{2}) 
\end{equation}
Importantly, the set of observations which are recorded at each time instant $t_{j}$ may differ. For each time instant $t_{j}$, we let $\calM_{j} \subseteq \calI$ represent the observations recorded at time $t_{j}$, with $M_{j} := \abs{\calM_{j}}$. 

The ultimate goal of the algorithm proposed by \citeauthor{Todescato-Efficient-2020} is to develop estimates of the \emph{entire} function $f(x, t)$ at any $x \in \calX$ and $t \in \bbR^{+}$ based solely on statistical inference on observations from the set of pre-specified locations $\calI \subset \calX$.

\begin{assumption}
\label{ass:todescato-kernel-sep}
    The kernel function $K$ representing the covariance of the univariate GP $f(x, t)$ can be separated into the product of a temporal and spatial kernel,
    \begin{equation}
    \label{eq:GP-kernel-sep}
        K((x, t), (x', t')) = K_{s}(x, x')\gamma(\tau), \quad \textrm{where}\: \tau := \abs{t - t'}
    \end{equation}
\end{assumption}

\begin{assumption}
\label{ass:todescato-psd-order}
    The power spectral density $S(\omega)$ (see \Cref{def:power-spectrum}) of $\gamma(\abs{t - t'})$ can be written as a rational function (with respect to $\omega$) of order $2r$. 
\end{assumption}

The separability requirement enforced by \Cref{ass:todescato-kernel-sep} restricts the class of kernel functions which can be used for describing the covariance structure of the Gaussian process. However, this assumption is not overly restrictive in practice, and the family of kernels meeting this criterion remains quite flexible. Similarly, \Cref{ass:todescato-psd-order} does not excessively confine the options for temporal kernels in the Gaussian process, as many commonly-used kernels such as the Mat{\'e}rn covariance family satisfy this requirement of the power spectrum \citep{GPML-Book-2006}.

In order to apply the Kalman filter, the algorithm proposed by \citeauthor{Todescato-Efficient-2020} relies on a discretized state-space representation of the dynamical process which guides the evolution of the function $f(x, t)$. In many applied settings, the dynamical system which is desired evolves according to a continuous-time process. However, the measurement model is based on discrete-time observations. To overcome this distinction, the authors introduce a discretized state-space representation of the underlying continuous-time dynamics which guide the `true' growth of the stochastic dynamical system of interest.

For any time instance $t_{j}$ where measurements are recorded, let $\tau_{j} := t_{j} - t_{j-1}$. Additionally, let the three matrices $F \in \bbR^{r \times r}, G \in \bbR^{r \times 1}$ and $H \in \bbR^{1 \times r}$ be the model matrices which collectively form a state-space representation for the spectral density $S(\omega)$, as described by \citeauthor{Carron-ML-Kalman-2016} and \Cref{thm:spectral-factor}.
\begin{definition}[Companion Form Matrices]
\label{def:companion-matrices}
    Suppose $S(\omega)$ is a rational power spectral decomposition for the covariance function $\gamma(\tau)$ with order $2r$. Let $S(\omega)$ have the spectral factorization (see \Cref{thm:spectral-factor}) $S(\omega) = W(\bmi\omega)W(-\bmi\omega)$, where 
    \begin{equation}
        W(\bmi\omega) := \frac{b_{r-1}(\bmi\omega)^{r-1} + b_{r-2}(\bmi\omega)^{r-2} + \cdots + b_{0}}{(\bmi\omega)^{r} + a_{r-1}(\bmi\omega)^{r-1} + \cdots + a_{0}},
    \end{equation}
    as in \Cref{eq:spectral-factor}. The companion form model matrices corresponding to the PSD $S(\omega)$ are the following:
    \begin{subequations}
        \label{eq:companion-model-matrices}
        \begin{align}
            F &= \begin{bmatrix}
                0 & 1 & 0 & \cdots & 0\\
                0 & 0 & 1 & \cdots & 0\\
                \vdots & \vdots & \vdots & \ddots & \vdots\\
                0 & 0 & 0 & \cdots & 1\\
                -a_{0} & -a_{1} & -a_{2} & \cdots & -a_{r-1}
            \end{bmatrix} \in \bbR^{r \times r} \label{eq:companion-model-matrix-F}\\
            G &= \begin{bmatrix}
                0 & 0 & 0 & \cdots & 1
            \end{bmatrix}^{\sfT} \in \bbR^{r \times 1} \label{eq:companion-model-matrix-G}\\
            H &= \begin{bmatrix}
                b_0 & b_{1} & b_{2} & \cdots & b_{r-1}
            \end{bmatrix} \in \bbR^{1 \times r} \label{eq:companion-model-matrix-H}
        \end{align}
    \end{subequations}
\end{definition}




Let $\overline{K}_{s} \in \bbR^{M \times M}$ be the Gram matrix for the spatial kernel $K_{s}$ applied to $\calI$. The spatiotemporal dynamics for the desired system are discretized by the following pair of equations, which guide the evolution from one observation time to the next:
\begin{subequations}
    \label{eq:todescato-discrete-evolution}
    \begin{align}
        \bfs(j+1) &= A(j)\bfs(j) + \bfw(j), \quad \bfw(j) \sim \calN(\mathbf{0}, Q(j)) \label{eq:todescato-discretized-state}\\
        \bfy(j) &= C(j)\bfs(j) + \bfv(j), \quad \bfv(j) \sim \calN(0, R(j)) \label{eq:todescato-discretized-measure}
    \end{align}
\end{subequations}
The Markovian dynamics which impact \Cref{eq:todescato-discretized-state,eq:todescato-discretized-measure} are parametrized by the following matrices. Formally, the four matrices given below are indexed by the observation time $t_{j}$, but the only manner in which they differ across observation times is through changes to the location set $\calM_{j} \subset \calI$. If the locations are always identical (i.e. $\calM_{1} = \dots = \calM_{J} = \calI$), then these matrices are fully constant. 
\begin{subequations}
    \label{eq:todescato-discretize-parameters}
    \begin{align}
        A(j) &= e^{\tau_{j}(\bfI_{M} \otimes F)} \in \bbR^{rM \times rM} \label{eq:todescato-discretize-Aj}\\
        C(j) &= I_{j}\overline{K}_{s}^{1/2}(\bfI_{M} \otimes H) \in \bbR^{M_{j} \times rM} \label{eq:todescato-discretize-Cj}\\
        Q(j) &= \bfI_{M} \otimes \left(\int_{0}^{\tau_{j}}\big(e^{\tau{F}}\big)GG^{\sfTsmall}\big(e^{\tau{F}}\big)\dx[\tau]\right) \in \bbR^{rM \times rM} \label{eq:todescato-discretize-Qj}\\
        R(j) &= \sigma^{2}\bfI_{M_{j}} \in \bbR^{M_{j} \times M_{j}} \label{eq:todescato-discretize-Rj}
    \end{align}
\end{subequations}
The matrix $I_{j}$ in \Cref{eq:todescato-discretize-Cj} is an indicator matrix $I_{j} \in \{0, 1\}^{M_{j} \times M}$ which selects for the subset of $x_{i} \in \calM_{j}$ from the entire location set $\calI$.

\begin{algorithm}
    \caption{Kalman-based GP regression on $\calI$ at query time $t \in \bbR^{+}$}
    \label{alg:todescato-kalman-GP}

    \begin{spacing}{1.1}
    \begin{algorithmic}[1]
        \Require 
            \Statex \hspace{\algorithmicindent} $(F,G,H)$: state-space representation of $S_r(\omega)$
            \Statex \hspace{\algorithmicindent} $\sigma^2$: measurement noise variance
            \Statex \hspace{\algorithmicindent} $\calI$: set of input locations
            \Statex \hspace{\algorithmicindent} $K_s(\cdot,\cdot)$: spatial kernel function
            \Statex \hspace{\algorithmicindent} $\gamma(\cdot)$: temporal kernel function
            \Statex \hspace{\algorithmicindent} $\bfSigma_0$: solution of Lyapunov equation $FX+XF^{\sfTsmall}+GG^{\sfTsmall}=0$

        \State Initialize $\widehat{\bfs}(0\mid -1) = \mathbf{0}$ and $\bfSigma(0\mid -1)=\bfI\otimes\bfSigma_0$.

        \Statex \texttt{Kalman estimates:}
        \While{$t_{j+1} \le t$}
            \State Compute $A(j), C(j), Q(j), R(j)$ using \Crefrange{eq:todescato-discretize-Aj}{eq:todescato-discretize-Rj}
            \State Apply the Kalman filter (\Cref{alg:kalman-filter}) to compute $\widehat{\bfs}(j+1 \mid j+1)$ and $\bfSigma(j+1 \mid j+1)$
            \State $\widehat{\mathbf{s}}(t) = \widehat{\mathbf{s}}(j+1\mid j+1)$
            \State $\bfSigma^{\bfs}(t)=\bfSigma(j+1\mid j+1)$
        \EndWhile 

        \Statex \texttt{Open-loop prediction:}
        \State $\tau = t - t_j$
        \State $\widehat{\bfs}(t) = e^{\tau(\bfI \otimes F)}\widehat{\bfs}(j \mid j)$
        \State $\bfSigma^{\bfs}(t) = \big(e^{\tau(\bfI \otimes F)}\big)\bfSigma(j|j)\big(e^{\tau(\bfI \otimes F)}\big)^{\sfTsmall}$

        \Statex \texttt{Process Estimate:}
        \State $\widehat{\bff}(t) = \overline{K}_{s}^{1/2}(\bfI \otimes H)\widehat{\bfs}(t)$
        \State $\bfSigma^{\bff}(t) = \overline{K}_{s}^{1/2}(\bfI \otimes H)\bfSigma^{\bfs}(t)(\bfI \otimes H)^{\sfTsmall}\overline{K}_{s}^{1/2}$
        \State \Return $\widehat{\bff}(t)$ and $\bfSigma^{\bff}(t)$
    \end{algorithmic}
    \end{spacing}
\end{algorithm}

\Cref{alg:todescato-kalman-GP} is a cleaned-up representation of the first algorithm presented by \citeauthor{Todescato-Efficient-2020} in their paper. Instead of presenting the algorithm as a continuous, non-terminating function which `loops' over all $t \in \bbR^{+}$, we have altered the formulation of the algorithm to explicitly demonstrate the Kalman filtering aspect, and the subsequent estimation of the change in the process post-filtering. 

Though it is not an inherent part of the algorithm, in practice, \Cref{alg:todescato-kalman-GP} should be treated as an \emph{online} algorithm, with the most recent Kalman-filtered state maintained in memory until a new set of observations are recorded and transmitted. Storing these terms in local memory will both allow for the estimation of current-time measurements based on a continuous-time stochastic dynamical equation, and it prevents the need to recompute previous states when subsequent measurements are transmitted and the updates to the current estimated states are desired.

As part of their paper, \citeauthor{Todescato-Efficient-2020} describe and prove a \emph{Dynamic Representer Theorem} \Cref{thm:dynamic-representer-theorem}; a form of the classic `Representer Theorem' for Gaussian processes (see \Cref{thm:representer-theorem}) which is amenable to the presumed spatiotemporal dynamics in their model. The key  practical result of their Dynamic Representer Theorem is that for any time $t$ beyond the final observation time $t_{J}$, the mean and variance estimates returned by \Cref{alg:todescato-kalman-GP} can be used to derive a minimum-variance estimator of the function $f(x, t)$ at \emph{any} location $x \in \calX$.

\begin{theorem}[Dynamic Representer Theorem]
\label{thm:dynamic-representer-theorem}
    Suppose that we have a Gaussian process $f: \calX \times \bbR^{+} \to \bbR$ with a spatiotemporal kernel $k(\cdot, \cdot)$ which satisfies \Cref{ass:todescato-kernel-sep,ass:todescato-psd-order}, and we have some observations recorded at times $t_{0}, \dots, t_{J}$. For any $x \in \calX$ and $t \ge t_{J}$, our estimate of $f(x, t)$ is the following:
    \begin{equation}
    \label{eq:dynamic-rep-mean}
        \widehat{f}(x, t) = \overline{K}_{s}(x, \calI)\overline{K}_{s}(\calI, \calI)\widehat{\bff}(t),
    \end{equation}
    and the variance of this estimator is given by 
    \begin{equation}
    \label{eq:dynamic-rep-var}
        \var{\widehat{f}(x, t) \mid \calD} = \gamma(0)\Big[\overline{K}_{s}(x, x) - \overline{K}_{s}(x, \calI)\overline{K}_{s}(\calI, \calI)^{-1}\big(\overline{K}_{s}(\calI, \calI) - \bfSigma^{\bff}(t)\big)\overline{K}_{s}(\calI, \calI)^{-1}\overline{K}_{s}(\calI, x)\Big]
    \end{equation}
\end{theorem}

The key concept behind the authors' proof of \Cref{thm:dynamic-representer-theorem} is that the assumptions regarding the evolution of the spatiotemporal dynamical system yield a strong conditional independence result. In particular, the value of $f(x, t)$ is \emph{conditionally independent} of the state of the dynamical system at any other time given the observations at time $t$. More formally, for any location $x \in \calX$ and measurement time $t \in \bbR$, we have that
\begin{equation}
\label{eq:todescato-cond-indep}
    f(x, t) \indep \bff(\tau) \mid \bff(t), \quad \tau \ne t
\end{equation}
The implication of \Cref{eq:todescato-cond-indep} is clear for any $x \in \calI_{j}$ when $t = t_{j}$, as we would therefore already have an observation corresponding to that location at that time. However, the key takeaway of this result is that the mean $\widehat{\bff}(t)$ and covariance $\bfSigma^{\bff}(t)$ given by the filtering distribution in \Cref{alg:todescato-kalman-GP} can be used as a posterior Gaussian process in estimating the value of $f(x, t)$ at \emph{any} location. As a corollary of this conditional independence result, the Representer Theorem \Cref{thm:representer-theorem} for Gaussian processes implies that $\widehat{f}(x, t)$ is a minimum-variance estimator for the value of $f(x, t)$.

In addition to defining and proving the Dynamic Representer Theorem, the authors describe the computational complexity of updating the estimated mean and variance in accordance with \Cref{alg:todescato-kalman-GP}. One iteration of the Kalman filtering algorithm to update the estimated mean and covariance from $\widehat{\bff}(t_{j-1})$ and $\bfSigma^{\bff}(t_{j-1})$ to $\widehat{\bff}(t_{j})$ and $\bfSigma^{\bff}(t_{j})$ has a computational complexity on the order of $\calO(rMM_{j} + rM_{j}^{3})$, with an additional cost of $\calO(MP)$ to compute predictions of the form $\{f(x_{p}, t_{j})\}_{i=1}^{P}$. Again, we reiterate here that this computation is made more efficient by treating \Cref{alg:todescato-kalman-GP} as an online algorithm, with some aspects of the model saved in persistent storage. For instance, extending the estimates to $P$ arbitrary locations requires computing $\overline{K}_{s}(\calI, \calI)^{-1}\overline{K}_{s}(\calI, x_{p})$ at each prediction location $x_{p}$. This matrix inversion and multiplication would typically be of the order $\calO(M^{2}P)$, but the matrix $\overline{K}_{s}(\calI, \calI)^{-1} \in \bbR^{M \times M}$ is constant across all iterations, so pre-computing and storing this matrix will allow for a significant practical speedup.

% \clearpage

\subsection{Adaptive Location Space}

In addition to the presentation of their primary Kalman-enhanced Gaussian process regression algorithm, \citeauthor{Todescato-Efficient-2020} also describe a second algorithm for 

\subsection{Multi-Output Dynamical System}
Now, we may focus on the adaptation of the Kalman-enhanced regression algorithm for a multi-output Gaussian process. In addition to the two previously-stated assumptions on the nature of the dynamical system, we must make one additional assumption regarding the observations in the dynamical model.
\begin{assumption}[Missingness of Vector-Valued Observations]
\label{ass:all-or-nothing}
    For any location $\bfx_{i} \in \calX$ and time $t_{k} \in \bbR^{+}$, the $D$ entries of the response vector $\bfy_{i,j} = f(\bfx_{i}, t_{j}) + \bfv_{i,j} \in \bbR^{D}$ will either all be recorded, or all be missing.
\end{assumption}
\Cref{ass:all-or-nothing} is a simplifying assumption which ensures that there are not any partially-recorded or partially-unrecorded vectors of information. Such an assumption was absent from the prior paper on this topic, as their algorithm focused solely on scalar-valued responses, so it would not be possible for the measurement $\bfy_{i}(t_{j})$ to be partially observed.

\begin{algorithm}
    \caption{Kalman-based MOGP regression on $\calI$ at query time $t \in \bbR^{+}$}
    \label{alg:kalman-MOGP-regression}

    \begin{algorithmic}[1]
        % \Require 
        %     \Statex \hspace{\algorithmicindent} $(F,G,H)$: state-space representation of $S_r(\omega)$
        %     \Statex \hspace{\algorithmicindent} $\sigma^2$: measurement noise variance
        %     \Statex \hspace{\algorithmicindent} $\calI$: set of input locations
        %     \Statex \hspace{\algorithmicindent} $K_s(\cdot,\cdot)$: spatial kernel function
        %     \Statex \hspace{\algorithmicindent} $\gamma(\cdot)$: temporal kernel function
        %     \Statex \hspace{\algorithmicindent} $\bfSigma_0$: solution of Lyapunov equation $FX+XF^{\sfTsmall}+GG^{\sfTsmall}=0$

        % \State Initialize $\widehat{\bfs}(0\mid -1) = \mathbf{0}$ and $\bfSigma(0\mid -1)=\bfI\otimes\bfSigma_0$.

        % \Statex \texttt{Kalman estimates:}
        % \While{$t_{j+1} \le t$}
        %     \State Compute $A(j), C(j), Q(j), R(j)$ using \Crefrange{eq:todescato-discretize-Aj}{eq:todescato-discretize-Rj}
        %     \State Apply the Kalman filter (\Cref{alg:kalman-filter}) to compute $\widehat{\bfs}(j+1 \mid j+1)$ and $\bfSigma(j+1 \mid j+1)$
        %     \State $\widehat{\mathbf{s}}(t) = \widehat{\mathbf{s}}(j+1\mid j+1)$
        %     \State $\bfSigma^{\bfs}(t)=\bfSigma(j+1\mid j+1)$
        % \EndWhile 

        % \Statex \texttt{Open-loop prediction:}
        % \State $\tau = t - t_j$
        % \State $\widehat{\bfs}(t) = e^{\tau(\bfI \otimes F)}\widehat{\bfs}(j \mid j)$
        % \State $\bfSigma^{\bfs}(t) = \big(e^{\tau(\bfI \otimes F)}\big)\bfSigma(j|j)\big(e^{\tau(\bfI \otimes F)}\big)^{\sfTsmall}$

        % \Statex \texttt{Process Estimate:}
        % \State $\widehat{\bff}(t) = \overline{K}_{s}^{1/2}(\bfI \otimes H)\widehat{\bfs}(t)$
        % \State $\bfSigma^{\bff}(t) = \overline{K}_{s}^{1/2}(\bfI \otimes H)\bfSigma^{\bfs}(t)(\bfI \otimes H)^{\sfTsmall}\overline{K}_{s}^{1/2}$
        \State \Return $\widehat{\bff}(t)$ and $\bfSigma^{\bff}(t)$
    \end{algorithmic}
\end{algorithm}

\section{Methods}
\label{sec:methods}

A major theoretical benefit of the algorithm proposed by \citeauthor{Todescato-Efficient-2020} -- and our multi-output extension in \Cref{alg:kalman-MOGP-regression} -- is that under the necessary assumptions presented above, these algorithms provide minimum variance estimators with a significant reduction in the required computational resources. In order to determine if this goal is met in a practical setting, it is necessary to assess the efficacy of the proposed algorithm based on testable, reproducible results.

To evaluate the performance of \Cref{alg:kalman-MOGP-regression} we compare the accuracy and speed of \Cref{alg:kalman-filter} to a `naive' multi-output Gaussian process regression algorithm. The code for comparing these algorithms are created in Python, and utilize the \texttt{GPyTorch} library \citep{GPyTorch-Library-2018} for the implementation of computational tasks involving Gaussian processes.

\subsection{Simulations}
The focus of our practical simulation study is to compare \Cref{alg:kalman-MOGP-regression} to a `traditional' multi-output Gaussian process regression method. 

To accomplish this, we generate synthetic data according to a multivariate Gaussian process model, and compare the predictive accuracy and overall runtime of these algorithms using a set of four multi-output Gaussian process kernels. The four kernels which are compared are described by the following equations:
\begin{subequations}
    \label{eq:simulation-kernels}
    \begin{align}
        \textbf{Independent:} \quad K((x, t, d), (x', t', d')) &= \gamma(\abs{t - t'}) K_{s}(x, x') \mathbf{1}_{\{d = d'\}}\sigma^{2}_{d}\\[6pt]
        \textbf{ICM:} \quad K((x, t, d), (x', t', d')) &= \gamma(\abs{t - t'}) K_{s}(x, x') \bfB(d, d')\\[6pt]
        \textbf{LMC:} \quad K((x, t, d), (x', t', d')) &= \gamma(\abs{t - t'}) \sum_{q=1}^{Q} K_{s}^{(q)}(x, x') \bfB^{(q)}(d, d')\\[6pt]
        \textbf{Temporal LMC:} \quad K((x, t, d), (x', t', d')) &= \sum_{q=1}^{Q} \gamma^{(q)}(\abs{t - t'}) K_{s}^{(q)}(x, x') \bfB^{(q)}(d, d')
    \end{align}
\end{subequations}
These four kernels describe similar, but distinct, covariance structures for a multi-output Gaussian process. The first two kernels are fully separable into the product of covariances across the temporal, spatial, and output dimensions, whereas the two LMC kernels are constructed via a sum of separable kernels, but they are not themselves separable. 
\begin{subequations}
    \label{eq:simulation-kernel-eqns}
    \begin{align}
        \textbf{Independent:} \quad K((x, t, d), (x', t', d')) &= \exp\!\Big(\!-\frac{\norm{x - x'}^{2}}{2}\Big)\exp(-\abs{t - t'})\mathbf{1}_{\{d = d'\}}\sigma^{2}_{d}\\[6pt]
        % \textbf{ICM:} \quad K((x, t, d), (x', t', d')) &= \gamma(\abs{t - t'}) K_{s}(x, x') \bfB(d, d')\\[6pt]
        % \textbf{LMC:} \quad K((x, t, d), (x', t', d')) &= \gamma(\abs{t - t'}) \sum_{q=1}^{Q} K_{s}^{(q)}(x, x') \bfB^{(q)}(d, d')\\[6pt]
        % \textbf{Temporal LMC:} \quad K((x, t, d), (x', t', d')) &= \sum_{q=1}^{Q} \gamma^{(q)}(\abs{t - t'}) K_{s}^{(q)}(x, x') \bfB^{(q)}(d, d')
    \end{align}
\end{subequations}



The simulation study was carried out on a computing server hosted by the Department of Statistical Sciences at the University of Toronto. To effectively run the simulations, we used \texttt{Python 3.12.5} \citep{python-3-12-5} with the \texttt{gpytorch} \citep{GPyTorch-Library-2018} library providing significant assistance.

% \subsection{Code Implementation}

% \subsection{Replication}

% \subsection{Extension}

\section{Results}
\label{sec:results}

\section{Discussion}
\label{sec:discussion}

\subsection{Conclusions}

In this paper, we focus on developing a method for deriving estimates of a spatiotemporal multi-output Gaussian process by incorporating a Kalman filter to sequentially update the mean and covariance functions.

% This work primarily follows the algorithm proposed by \citeauthor{Todescato-Efficient-2020} to couple Gaussian process regression with Kalman filtering methods for scalar outcome variables.

\subsection{Limitations}

\subsection{Future Work}

One potential area of continued algorithmic research related to this project is to examine the use of smoothing distributions such as the Rauch-Tung-Striebel (RTSS) smoother \citep{Rauch-RTSS-1965,Bayesian-Filter-Book-2023} or modified Bryson-Frazier smoother \citep{Colemont-MBF-Smoothing-2026}. The primary use of these smoothing distributions is to develop estimates for the latent states $f(x_{i}, t_{j})$ at \emph{any} observation time $t_{k}$ instead of solely performing inference on the latent states corresponding to the latest-observed measurements. The RTSS has been shown \citet{Hartikainen-Smoothing-GP-2010,Sarkka-Spatial-Kalman-2012} to be closely related to temporal Gaussian processes, which can enable the use of smoothers for the multi-output Gaussian process problems considered in this work.

Another related task of future interest is to examine the possibility of working with dynamical systems whose transformations are governed by unknown parameters. In this work, we have focused on dynamical systems where the mechanics underpinning the evolution of the system over time are fully known. However, this is a reductive assumption, which may be improved upon through incorporating unknown system parameters into the state transition model, instead of treating the evolution equations as fully specified and completely understood. For example, a parametric or semi-parametric kernel function may have hyperparameters which can be inferred jointly with the latent states in the dynamical system. Developments in this research area may be particularly relevant in real-world scientific research applications in fields such as epidemiology, environmental science, and control systems engineering, as real-world dynamical systems are often governed by time-varying behaviour, and we may lack complete knowledge of the corresponding mechanics driving the evolution of the system.

% \bigskip
% \bigskip


% References
\newpage

\bibliographystyle{apalike}
\bibliography{references}


% Appendix
\newpage
\appendix
\section*{Appendix}

\section{AI Usage Statement}
Two large language models -- Claude Opus 5.5 (Anthropic) and GPT-5.6 Luna (OpenAI) -- were used in the creation of this report.

Claude Opus 5.5 was employed for the creation and testing of Python code used for the simulations and replication of results from the prior paper. Claude Opus 5.5 was additionally used for assistance with some mathematical derivations and proofs within the written report. GPT-5.6 Luna was used for \LaTeX typesetting and formatting. Both large language models were used in searching for relevant background literature pertaining to this project.

All of the content in this document has been verified by the author, and is the sole responsibility of the author of this document. All AI-generated work in the creation of this project has been thoroughly reviewed and revised upon by the author, who is solely responsible for the content therein.

\section{Code} 
\label{sec:code} 
For the sake of brevity, the code utilized in the creation of this project is not directly present in the appendix. All of the code used for this project is available in a public GitHub repository, located at \url{https://github.com/hotz1/Dynamical-Multitask-Prediction}. The \texttt{README.md} file in the root directory of the GitHub repository contains further information regarding these files.

\section{Definitions and Technical Details}
\label{sec:definitions}

\begin{definition}[Reproducing Kernel Hilbert Space (RKHS)]
\label{def:rkhs}
    Let $\calX$ be an indexing set, and let $\calH$ be a Hilbert space of real-valued functions $f: \calX \to \bbR$ equipped with the inner product $\innerprod[\calH]{\cdot}{\cdot}$ and induced norm $\norm[\calH]{\cdot}$. 
    
    For every $x \in \calX$, define the evaluation functional $L_{x}: \calH \to \bbR$ such that $L_{x}[f] = f(x)$ for all $f \in \calH$. $\calH$ is an RKHS over $\calX$ if for all $x \in \calX$, the evaluation functional $L_{x}[f] := f(x)$ is a bounded and continuous functional on $\calH$. 
\end{definition}

\begin{definition}[Reproducing Kernel]
\label{def:reproducing-kernel}
    Let $\calH$ be an RKHS over $\calX$ with $\innerprod[\calH]{\cdot}{\cdot}$ and norm $\norm[\calH]{\cdot}$. The function $k: \calX \times \calX \to \bbR$ is the reproducing kernel for $\calH$ if it satisfies the reproducing property:
    \begin{equation}
        f(x) = \innerprod[\calH]{f}{k(x, \cdot)}, \quad \forall x \in \calX
    \end{equation}
\end{definition}

\begin{theorem}[Representer Theorem]
\label{thm:representer-theorem}
    Suppose that we have some dataset $\calD := \{(x_{i}, y_{i})\}_{i=1}^{n}$.
    \begin{equation}
        J[f] := \frac{\lambda}{2}\norm[\calH]{f}^{2} + \sum_{i=1}^{n}L\big(y_{i}, f(x_{i})\big),
    \end{equation}
    where $\lambda$ is a regularization parameter, and $L$ is some loss function. Let $\calH_{k}$ be an RKHS with reproducing kernel $k$. The Representer Theorem states that the minimizer $f^{*}(x)$ satisfies
    \begin{equation}
        f^{*}(x) = \argmin_{f \in \calH_{k}}J[f] = \sum_{i=1}^{n}\alpha_{i}k(x, x_{i})
    \end{equation}
\end{theorem}

\begin{definition}[Power Spectral Density]
\label{def:power-spectrum}
    The power spectral density (PSD) of a univariate real-valued function $z(\tau)$ is given by
    \begin{equation}
    \label{eq:PSD-definition}
        S(\omega) = \int z(\tau)e^{-2\pi{i}\omega\tau}\dx[\tau]
    \end{equation}
\end{definition}

\begin{theorem}[Wiener-Khinchin Theorem]
\label{thm:wiener-khinchin}
    Suppose that $S(\omega)$ is the power spectral density of a univariate real-valued function $z(\tau)$, where $z(\tau)$ is the covariance of a stationary Gaussian process with $\tau := t - t'$. The PSD $S(\omega)$ is equal to the Fourier transform of the covariance $z$, given by 
    \begin{equation}
        S(\omega) := \calF[z(\tau)](\omega)
    \end{equation}
\end{theorem}

\begin{theorem}[Spectral Factorization]
\label{thm:spectral-factor}
    Suppose that $S(\omega)$ is the power spectral density of a univariate real-valued function $z(\tau)$, and $S(\omega)$ is a rational function of order $2r$. The power spectral density $S(\omega)$ can be decomposed into $S(\omega) = W(\bmi\omega)W(-\bmi\omega)$, where $W(\cdot)$ is described by the following:
    \begin{equation}
    \label{eq:spectral-factor}
        W(\bmi\omega) := \frac{b_{r-1}(\bmi\omega)^{r-1} + b_{r-2}(\bmi\omega)^{r-2} + \cdots + b_{0}}{(\bmi\omega)^{r} + a_{r-1}(\bmi\omega)^{r-1} + \cdots + a_{0}}
    \end{equation}
\end{theorem}
\end{document}
